\section{Reconciliation with the Budget Lab}\label{sec:reconciliation}

The two models measure different totals: theirs is federal individual income
tax plus payroll plus a corporate wedge, with no states and no non-credit
transfers; this paper's household total spans federal and state and nets out
the whole transfer system. Setting the headline numbers side by side without
a bridge would be wrong, so this section builds one, using only this paper's
run output and objects they publish. Their repository commits the complete
published result grid --- revenue by instrument for all eighteen cells ---
which makes the comparison possible at the cell level rather than against
two text-stated totals.

\textbf{Basis.} Their microsimulation total is individual income tax plus
payroll minus tax-credit outlays; the matching quantity here is federal
income tax net of refundable credits, plus payroll (both halves and
self-employment). Their corporate wedge reduces to the capital growth rate
times CBO's 2030 corporate anchor of \$486B --- in their formula
$\Delta R_{\mathrm{CIT}} = X \cdot (R_{\mathrm{CIT}}/Y^0_K)$ with
$X = Y^0_K\,g$, the capital base cancels --- and all eighteen committed
values equal $g \times \$486\mathrm{B}$ to within \$0.15B. Adding their wedge
to this paper's federal microsimulation total yields a like-for-like
``bridged'' headline. The wedge is not PolicyEngine modelling corporate tax,
and it inherits every constant-rate assumption they list.

\begin{table}[t]
\centering
\tabBridge
\caption{The bridged headline cell (Rapid~/~expansive, their central
published scenario), \$B, 2030. ``IIT (net)'' is individual income tax net
of refundable credits; the corporate wedge is theirs in both rows.}
\label{tab:bridge}
\end{table}

\textbf{Where the models agree.} Payroll responses agree closely in every
Rapid cell (Table~\ref{tab:cells}): ours run
\genPayrollTripletOurs{} against theirs \genPayrollTripletTheirs{} across
compressive~/~proportional~/~expansive --- two independent models putting
similar wage distributions against the same taxable maximum, including the
sign reversal under compression. The tilt ratios also line up once stated on
the same basis: on their federal-plus-corporate accounting, shares-fixed
revenue exceeds as-forecast revenue by
\genTiltOursMatchedSlow$\times$~/~\genTiltOursMatchedMod$\times$~/~\genTiltOursMatchedRapid$\times$
here against
\genTiltTheirsSlow$\times$~/~\genTiltTheirsMod$\times$~/~\genTiltTheirsRapid$\times$
in their grid for Slow~/~Moderate~/~Rapid --- agreement on the paper's
central mechanism, with the same ordering. (The household-total ratios here
are larger --- \genTiltOursHouseholdRapid$\times$ for Rapid --- because
states and transfers amplify what the tilt costs; that is a coverage
statement, not a disagreement.) Their most extreme published comparison also
reproduces directionally: for Slow with compressive inequality, their grid
puts revenue \genSlowCompTheirsLowerPct\% below the shares-and-distribution-%
unchanged counterfactual (their text states 82\%); this paper gets
\genSlowCompOursLowerPct\%.

\begin{table}[t]
\centering
\tabPayrollCells
\caption{Cell-level comparison on the matched federal basis, Rapid variants,
\$B, 2030.}
\label{tab:cells}
\end{table}

\textbf{Where they diverge, and why.} The bridged headline is
\genBridgeOursTotal B against their \genBridgeTheirsTotal B ---
\genBridgeRatio$\times$, or \genBridgeOursCboPct\% against
\genBridgeTheirsCboPct\% of CBO 2030 revenue. The gap concentrates in the
income-tax response to the capital shock: on the proportional cell, where
the wage-spread channel is neutral and the income-tax response is almost
entirely the capital shock, this paper's response is
\genIITRatioProp$\times$ theirs (\genIITRatio$\times$ on the expansive
bridge cell). Three model differences point the same direction. First, the
realized-capital base: their methodology's remark that the effective
corporate parameter ``falls out to roughly one'' implies a baseline realized
taxable capital base near \$4.6T, against \$\genModelledCapitalShareT T
here, roughly 14\% larger. Second, allocation: they distribute the
incremental capital flow by SCF-imputed total wealth; this paper distributes
by existing realized capital income, which they note is the more
top-concentrated choice --- and top-concentrated dollars face higher
marginal rates. Third, retirement routing: taxable pension and IRA
distributions here scale with the capital shock and are taxed at ordinary
rates. The capital-scope sensitivity of Section~\ref{sec:scope} shows that
last channel alone is worth about \$105B in Rapid.
